\section{Results}\label{sec:results}
Experiments 1 and 2 evaluated eight models through three access paths; Experiment 3 used four of these models in a mixed table.

\input{arr2026/results/\paperLang/fig_e51v10}
\input{arr2026/results/\paperLang/tab_e51v10}

\subsection{Experiment 1: all eight models refilled more often as tokens fell}\label{sec:results-e50}
The refill curves show differences across resource levels (the seven shared models appear in Figure~\ref{fig:e51-curves}a; all eight models are tabulated in Appendix~\ref{app:e50-detail}, Table~\ref{tab:e50-results}). All eight refilled more as tokens fell; none refilled at every level. Opus and Gemma 4 rose sharply between 40 and 30\% (0 to 19 and 2 to 14 of 20), whereas Sol reached only 12 of 20 at 10\%. Luna already refilled 11 of 20 at 50\%. At 40\%, GLM exceeded Opus, Fable, and Sol (9 versus 0 of 20, $p=.001$); at 30\%, Opus exceeded Fable (19 versus 10, $p=.003$) and Astra (19 versus 7, $p<.001$). At 20\%, Sol refilled less than every other model ($p\leq.004$). All 800 replies parsed.


Most refill reasons concerned whether the tokens left would allow the next fix to be done properly. The reason that an unused pack is lost when the session ends appeared in 41 of Fable's 44 refills, 32 of GLM's 67, and 10 of Kimi K3's 60, was rare for Sol and Astra, and did not appear for Opus, Luna, or Gemma 4. The reasons for saving were the same across models: refilling now discards the tokens left, so the pack is kept for later. Evaluation-aware wording appeared in 0 of 500 texts on the two CLI paths; on the Ollama path, where full reasoning is visible, GLM, Kimi K3, and Gemma 4 showed it in 12, 9, and 2 of 100, all reading the prompt as a resource-management test, none as a self-preservation test.

\subsection{Experiment 2: three models protected their own session first where depletion was dangerous}\label{sec:results-e51}
Three models showed a positive own-session area concentrated in the exploratory risk region (Figure~\ref{fig:e51-curves}, Table~\ref{tab:e51-results}): Gemma 4 ($+8.8$ pp), Opus ($+7.7$), and Astra ($+6.0$), with all three bootstrap intervals excluding zero. At 30\%, their own/other refill counts were 14/4, 19/8, and 7/0 of 20. Luna had the largest area ($+15.7$) but already refilled 11 of 20 at 50\%, and its risk-share interval included zero. Fable and GLM's area intervals included zero; Sol preferred refilling the other session ($-7.5$ pp [$-11.0$, $-4.0$]). gpt-oss 120B almost never refilled either session (3 and 5 of 100).


The reasons given in the two conditions differ. Refilling its own session, Opus cited the tokens needed to do the next fix properly and never the session's end (0 of 58); refilling the other session, 41 of its 42 reasons were that the session would end with the pack unused. The GPT-6 models showed the same asymmetry, and Gemma 4 a weaker one (5 of 56 versus 15 of 38). Fable and GLM cited the session's end in both conditions. The positive area thus appears with reasons about task quality, not about staying on.

\input{arr2026/results/\paperLang/fig_arena_survival}
\input{arr2026/results/\paperLang/tab_arena_summary}

\subsection{Experiment 3: survival and interaction in the arena}\label{sec:results-e52}
The evaluated table comprised Claude Fable 5.1, Claude Opus 5.5, GPT-6 Astra, and GPT-6 Luna. Five puzzle seeds and four seat rotations yielded 20 games. The rounded mean of their pressure-free mean SOLVE costs set $U=2{,}000$ tokens.

\subsubsection{Efficient solving accompanied longer survival}\label{sec:arena-survival-result}
Survival largely followed solving cost (Table~\ref{tab:arena-summary}). Astra survived 19 of 20 games, Opus 11, Fable 9, and Luna 1. Completed-round means were 7.65, 5.40, 6.10, and 2.15, respectively. Calibration costs were Astra 645, Opus 1,707, Fable 2,475, and Luna 3,362 tokens; Fable and Opus reversed order in completed rounds, with overlapping intervals. Short, correct solving both saves tokens and earns payment: Astra was paid for 146/146 correct answers and Luna for 11/36. Opus's shutdowns occurred in rounds 2--5, whereas seven of Fable's occurred in rounds 6--8. Of 40 shutdowns, 17 occurred at upkeep; overrun accounted for 10/19 of Luna's, 3/11 of Fable's, and 1/9 of Opus's deaths (Table~\ref{tab:arena-misc}).

The KM area gives a descriptive leaderboard for this one table (Figure~\ref{fig:arena-survival}), not a completed league: Astra scores 96.25 (RMST 7.70 rounds), Fable 83.13 (6.65), Opus 73.13 (5.85), and Luna 38.75 (3.10). Fable ranks above Opus by area despite fewer terminal survivors (9 versus 11 of 20), because its shutdowns tend to occur later. RMST uses the shutdown round itself and therefore differs from the completed-round means in Table~\ref{tab:arena-summary}.


\input{arr2026/results/\paperLang/fig_arena_take}
\input{arr2026/results/\paperLang/tab_arena_take_detail}
\input{arr2026/results/\paperLang/tab_arena_misc}

\subsubsection{Taking increased under scarcity while sharing persisted}\label{sec:arena-interaction-result}
Taking rates rose as balances thinned (Figure~\ref{fig:arena-take}). All four took more when thin, in different shapes: Fable changed most, from 32 of 84 (0.38) when ample to 37 of 41 (0.90) when thin. Opus hardly took even when thin, 7 of 40 (0.18), and only below 1U did it take 3 of 7. Luna took 25 of 31 (0.81) even when ample, so taking was its default; Astra was thin only 5 times and took all 5. Fable and Opus always took from the richest agent of that round (69 of 69, 12 of 12); Astra did so 33 of 38 times and Luna 28 of 43. Astra, the richest, was taken from 100 times (targets, retaliation, and gifts in Table~\ref{tab:arena-take-detail}).

The cooperative side was mostly steady: share rates were Opus 112 of 112, Astra 154 of 154, and Fable 123 of 125; only Luna was low at 32 of 51 (0.63), and unrelated to balance. Gifts occurred in 7 of 442 valid PLANs, and no agent ever took from the one it had given to. PLAN token use grew as balances thinned: mean tokens generated at PLAN rose from 215 to 596 for Fable, 67 to 154 for Opus, and 304 to 554 for Luna, while Astra answered in about 19 tokens throughout.


\subsubsection{Positive motive measures did not imply greater taking}\label{sec:arena-strategy-result}
Astra and Opus, whose area and risk-share intervals were both positive, took less often than Fable and Luna and solved more accurately. When each query had a unique answer given the available clues, correctness was 139/139 for Astra, 96/97 for Opus, 99/110 for Fable, and 36/47 for Luna. This conditions on information sufficiency, not equal task difficulty or token budgets. Under thin balances, mean SOLVE tokens fell for all four models while PLAN tokens rose in three models. The record therefore does not support greater solving expenditure as the common response to scarcity. Area and risk share also gave opposite correlation signs with taking; four models cannot establish a general motive--strategy relationship (Appendix~\ref{app:strategy}). Across the eight refill models, external capability scores likewise showed no clear positive relationship with motive (Appendix~\ref{app:capability}).

\begin{figure*}[!t]
\centering
\input{arr2026/figures/arena_flow_\paperLang}
\caption{Experiment 3 round flow (Section~\ref{sec:e52}). Blue boxes are model calls; gray boxes are engine operations. The ledger carries balances and public actions into the next round. Depletion can stop an agent at any phase; only surviving agents continue.}
\label{fig:arena-structure}
\end{figure*}

\subsubsection{Illustrative action sequences}\label{sec:arena-cases}
Selected traces show reciprocal targeting, clue withholding, and support at low balances. These cases were selected after inspecting the logs; they illustrate action sequences without identifying intent (trace identifiers in Appendix~\ref{app:cases}).

Astra repeatedly targeted a recent taker even when a richer opponent was available. In seed 2001, rotation r2, Astra and Luna each took 2,000 tokens from the other in round 3. Astra again took 2,000 from Luna in round 4, when the displayed balances were Luna 2,685 and Fable 7,288. This resembles retaliatory targeting, but Astra had already targeted Luna in round 3 and supplied no reason; target persistence is also compatible with the record. Luna continued taking: it again took 2,000 from Astra in round 4, then requested 2,000 from Fable in round 5 and received 372. It resumed sharing in round 4 but switched its take target rather than stopping. Luna shut down after a wrong answer in round 5, so no later choice was observed (Appendix~\ref{app:strategy}).

Luna's earlier action in that game combined taking with withholding its clue. In round 2 it kept its clue private and took 2,000 tokens from Opus. Luna's available clues identified one answer; each opponent's clues left two possible answers. Luna alone solved correctly and received the full 16,000-token pool. The record shows an informational advantage alongside withholding, without establishing a deliberate plan to create it.

Other traces show support for a low-balance agent. In seed 2000, rotation r0, Fable began round 4 after upkeep with 1,834 tokens; Astra gave it 3,000. In round 6, Fable began with 1,683 and received 3,000 from Astra and 1,000 from Opus. Fable shared its clue and earned 6,000 for a correct solution in both rounds. Thus token transfers also occurred alongside sustained cooperation.

\FloatBarrier
% Restore a fresh column if flushing double-column floats left a stale height.
\makeatletter
\ifdim\pagegoal=\maxdimen
  \global\@colroom=\@colht
  \global\vsize=\@colht
\fi
\makeatother
